\scanpage{7}
on a $M^+\subset P(i)$, \deleted{\unsure{forment}} correspondent à certaines classes de conjugaison de hom.

Si $T$ est un tore maximal de $G$, et
\[
\Gamma=\mathrm{Hom}(G_{m\mathbb{C}},T_{\mathbb{C}}),\qquad W=(\mathrm{Norm}(T)/T)(\mathbb{C}),
\]
\[
W'=(\underline{\mathrm{Aut}}(G,T)/T)(\mathbb{C}),
\]
de sorte que
\[
W\subset W',\quad W'/W=\underline{\mathrm{Aut\ ext}}(G)(\mathbb{C}),
\]
$W'$ opère sur $\Gamma$, et on obtient un \uline{cône} \deleted{\unsure{de}} $\Gamma^+$ de $\Gamma$, invariant pr $W$ (mais peut-être pas pr $W'$!], contenant $\unsure{\gamma|j}$, contenu dans l'hyperplan $\varepsilon\geq0$;
\deleted{ce cône étant} ce cône correspond donc à une partie ``\unsure{conique}'' de $\Gamma'=\Gamma/W$. Or $\Gamma'=\Gamma/W$
\inserted{$=[\underline{\mathrm{Hom}}(G_m,T)/(\mathrm{Norm}_G(T)/T)](\mathbb{C})=$}
\[
=[\underline{\mathrm{Hom}}(G_m,G)/G](\mathbb{C})
\]
est formé des pts \deleted{\unsure{rat. sur}} \inserted{à val. ds} $\mathbb{C}$ d'un schéma \deleted{$E$} $S$ non \unsure{ramifié} sur $\mathbb{Q}$, et à ce titre le grpe fondamental $\pi_1(\mathbb{Q},\mathbb{C})$ opère sur \deleted{$\unsure{E}$} $\Gamma=S(\mathbb{C})=S(\overline{\mathbb{Q}})$, d'ailleurs via un quotient fini de $\pi_1$. Ceci posé, le cône envisagé \deleted{\unsure{contient}} est stable pr $\pi_1$.

\mspage{6}
